\documentclass[12pt]{article}
\usepackage[spanish]{babel}
\usepackage{natbib}
\usepackage{url}
\usepackage{float}
\usepackage[utf8]{inputenc}
\usepackage{amsmath}
\usepackage{graphicx}
\usepackage{fancyhdr}
\usepackage{longtable}
\usepackage{vmargin}
\usepackage{listings}
\usepackage{hyperref}
\usepackage{multirow}
\setmarginsrb{3 cm}{2.5 cm}{3 cm}{2.5 cm}{1 cm}{1.5 cm}{1 cm}{1.5 cm}

\title{Tarea - 03}
\date{\today}

\makeatletter
\let\thetitle\@title
\let\theauthor\@author
\let\thedate\@date
\makeatother

\pagestyle{fancy}
\fancyhf{}
\rhead{\theauthor}
\lhead{\thetitle}
\cfoot{\thepage}

\addto\captionsspanish{
  \renewcommand{\contentsname}%
    {Tabla de contenido}%
}

\begin{document}
    \pagestyle{fancy}
    \fancyhf{}

    \lhead{\begin{picture}(0,0) \put(0,0){\includegraphics[width=30mm]{images/Logo2.png}} \end{picture}}
    \renewcommand{\headrulewidth}{0.7pt}
    \fancyhead[R]{Procesamiento de lenguaje natural - Tarea 02}
\fancyfoot[R]{\thepage}

\begin{titlepage}
	\centering
    \includegraphics[scale = 0.45]{images/Logo.png}\\[0.5 cm]	% University Logo
    \textsc{\large Universidad de los Andes\\
        \vspace{0.2cm} 
        Facultad de ingeniería\\
        \vspace{0.3cm} 
        Tarea 03}\\[2.0 cm]	% University Name
	\textsc{\Large Procesamiento de lenguaje natural}\\[0.5 cm]
	% Course Code
	\rule{\linewidth}{0.2 mm} \\[0.4 cm]
	{ \LARGE \bfseries \thetitle}\\
	\rule{\linewidth}{0.2 mm} \\[1.5 cm]
	
	\large
			\emph{Presentado por:} \\
			Juan David García Hernández\\
			Nicolás Rocha Pacheco\\
			César Daniel Garrido Urbano\\
	\vfill
	\large
			\emph{Presentado a:}\\
			Rubén Francisco Manrique Pirmanrique\\
\end{titlepage}

\thispagestyle{empty}
\tableofcontents
\pagebreak

\setcounter{page}{1}

\newpage
\bibliographystyle{unsrt}
\bibliography{biblist.bib}

\end{document}